
\subsection{规范不变性与范畴 $\mathfrak{HA}_\infty(C_1,C_0)$}
\label{subsec:5-2-4}

在证明定理 5.2.35 之前，
先应用该定理说明：由 $b_1,b_0$ 形变而得的
滤过 $A_\infty$ 双模
$(D,{}^{b_1}\mathfrak n^{b_0})$ 的同伦类，
不依赖于在各自规范等价类中所选的代表 $b_i$。
特别地，弗洛尔上同调在 $b_i$ 的规范等价下不变
\liftIndex{I188}。
$\widehat{\mathcal M}(C)$ 上规范等价关系的定义见第 4.3 节。

\begin{proposition}{5.2.37}\label{prop:5-2-37}
\begin{enumerate}[label=\textup{(\arabic*)},leftmargin=2.8em]
\item 设 $(D,\mathfrak n)$ 为
$(C_1,C_0)$ 上的滤过 $A_\infty$ 双模，
且 $b_1,b'_1\in\widehat{\mathcal M}(C_1)$、
$b_0,b'_0\in\widehat{\mathcal M}(C_0)$。
若 $b_1\sim b'_1$ 且 $b_0\sim b'_0$，
则 $(D,{}^{b_1}\mathfrak n^{b_0})$
与 $(D,{}^{b'_1}\mathfrak n^{b'_0})$
严格同伦等价。
\item 假设 $C_i$（$i=0,1$）与 $D$ 含幺
\translatornote{原书实体页 297 在此还列出
$C'_i$，但命题 (2) 并未定义或使用它；
本项只涉及 $C_i$ 上的定界余链
及其双模 $D$，故不把未定义的
$C'_i$ 带入假设。见疑误 A-0276。}，
并取
$b_1,b'_1\in\widehat{\mathcal M}_{\mathrm{weak}}(C_1)$、
$b_0,b'_0\in\widehat{\mathcal M}_{\mathrm{weak}}(C_0)$。
若 $b_1\sim b'_1$、$b_0\sim b'_0$，
且
$\mathfrak{PO}_1(b_1)=\mathfrak{PO}_0(b_0)$
\translatornote{原书实体页 297
仍像命题 4.3.25 一样省去两端势函数的下标；
两端元素分属 $C_1,C_0$，
这里按相应定义域标明
$\mathfrak{PO}_1,\mathfrak{PO}_0$。
见疑误 A-0279。}，
则 $(D,{}^{b_1}\mathfrak n^{b_0})$
与 $(D,{}^{b'_1}\mathfrak n^{b'_0})$
\emph{su-同伦等价}。
参见定义 5.2.34 与 4.2.42。
同伦含幺版本也成立。
\item 设 $(L^{(0)},L^{(1)})$
为相对 spin 的拉格朗日子流形对，
并对 $i=0,1$ 取
$(\mathfrak b,b_i)$、
$(\mathfrak b',b'_i)
\in\widehat{\mathcal M}_{\mathrm{weak,def}}
(L^{(i)})$
\translatornote{原书实体页 297
此处写无帽记号
$\mathcal M_{\mathrm{weak,def}}$；
但其后的 $\sim$ 是定义--推论 4.6.50
施于未取商的
$\widehat{\mathcal M}_{\mathrm{weak,def}}$
上的规范等价关系，
推论 5.2.40 (3) 亦使用带帽记号。
见疑误 A-0277。}。
假设
$\mathfrak{PO}_{L^{(0)}}(\mathfrak b,b_0)
=\mathfrak{PO}_{L^{(1)}}(\mathfrak b,b_1)$。
若
$(\mathfrak b,b_i)\sim
(\mathfrak b',b'_i)$
（$i=0,1$），
则含幺的滤过 $A_\infty$ 双模
\translatornote{原书实体页 297 在此写
\emph{fibered} $A_\infty$ bimodule，
但本项讨论的是诺维科夫环上的形变滤过
$A_\infty$ 双模，且“fibered”并未在此定义；
按本章既有对象译作“滤过”，
见疑误 A-0278。}
\[
\bigl(C(L^{(1)},L^{(0)};\Lambda_{0,\mathrm{nov}}),
{}^{(\mathfrak b,b_1)}
\mathfrak n^{+(\mathfrak b,b_0)}\bigr)
\]
与
\[
\bigl(C(L^{(1)},L^{(0)};\Lambda_{0,\mathrm{nov}}),
{}^{(\mathfrak b',b'_1)}
\mathfrak n^{+(\mathfrak b',b'_0)}\bigr)
\]
\emph{shu-同伦等价}。
参见引理--定义 3.8.75。
\end{enumerate}
在 (1)--(3) 中，
均可选取模 $\Lambda^+_{0,\mathrm{nov}}$
同余于恒等同态的同伦等价。
\translatornote{原书实体页 297
在 “congruent to the identity” 之后
写作 “module $\Lambda^+_{0,\mathrm{nov}}$”；
结合第 5.2.2 小节同类断言，
这里应为 “modulo $\Lambda^+_{0,\mathrm{nov}}$”。
见疑误 A-0280。}
\end{proposition}

\begin{proof}
令
$\widetilde b_i
\in\widehat{\mathcal M}(\mathfrak C_i)$
为从 $b_i$ 到 $b'_i$ 的同伦。
取 $[0,1]\times D$ 的模型
$\mathfrak D$，
它是
$(\mathfrak C_1,\mathfrak C_0)$-双模。
由于
\[
\mathrm{Eval}_{s=0}\widetilde b_i=b_i,
\qquad
\mathrm{Eval}_{s=1}\widetilde b_i=b'_i,
\]
由定义--引理 5.2.18 可知，
两个端点评价分别诱导
$A_\infty$ 双模同态
\[
\begin{aligned}
\mathrm{Eval}_{s=0}:
(\mathfrak D,
{}^{\widetilde b_1}\mathfrak n^{\widetilde b_0})
&\longrightarrow
(D,{}^{b_1}\mathfrak n^{b_0}),\\
\mathrm{Eval}_{s=1}:
(\mathfrak D,
{}^{\widetilde b_1}\mathfrak n^{\widetilde b_0})
&\longrightarrow
(D,{}^{b'_1}\mathfrak n^{b'_0}).
\end{aligned}
\]
由定理 5.2.35，
二者都是严格同伦等价，
故 (1) 得证。
利用引理 4.3.23，
以同样方式证明 (2)。

为证 (3)，先取定理 4.6.47 的同态
\[
\mathfrak f_i:
(C_i,\mathfrak m^{\mathfrak b})
\longrightarrow
(C'_i,\mathfrak m^{\mathfrak b'}).
\]
由引理 5.2.16、5.2.17，
有严格同伦等价
\[
\mathfrak f_i^{b_i}:
(C_i,\mathfrak m^{(\mathfrak b,b_i)})
\longrightarrow
(C'_i,\mathfrak m^{
(\mathfrak b',\,\mathfrak f_{i*}b_i)}).
\]
因此，由引理 5.2.19，
在
$(\mathfrak f_1^{b_1},\mathfrak f_0^{b_0})$
之上有严格同伦等价
\[
\bigl(D,
{}^{(\mathfrak b,b_1)}
\mathfrak n^{(\mathfrak b,b_0)}\bigr)
\sim
\bigl(D,
{}^{(\mathfrak b',\,\mathfrak f_{1*}b_1)}
\mathfrak n^{
(\mathfrak b',\,\mathfrak f_{0*}b_0)}\bigr).
\]
另一方面，
由定义--推论 4.6.50，
$(\mathfrak b,b_i)\sim
(\mathfrak b',b'_i)$
蕴含
$\mathfrak f_{i*}b_i\sim b'_i$
于
$\widehat{\mathcal M}_{\mathrm{weak}}
(C'_i,\mathfrak m^{\mathfrak b'})$。
于是类似上述论证给出
\[
\bigl(D,
{}^{(\mathfrak b',\,\mathfrak f_{1*}b_1)}
\mathfrak n^{
(\mathfrak b',\,\mathfrak f_{0*}b_0)}\bigr)
\sim
\bigl(D,
{}^{(\mathfrak b',b'_1)}
\mathfrak n^{(\mathfrak b',b'_0)}\bigr).
\]
故 (3) 得证。
命题最后一句由构造直接可见。
\end{proof}

\begin{proposition}{5.2.38}\label{prop:5-2-38}
\begin{enumerate}[label=\textup{(\arabic*)},leftmargin=2.8em]
\item 设 $\varphi,\psi:D\to D'$
为位于严格滤过 $A_\infty$ 同态之上的
滤过 $A_\infty$ 双模同态。
若 $\varphi\sim_s\psi$，则
\[
\varphi_*=\psi_*:
H^*(D,\mathfrak n_{0,0})
\longrightarrow
H^*(D',\mathfrak n'_{0,0}).
\]
\item 位于严格滤过 $A_\infty$ 代数的
严格同伦等价之上的滤过
$A_\infty$ 双模同伦等价，
在 $\mathfrak n_{0,0}$ 上同调上诱导同构。
\end{enumerate}
\end{proposition}

\begin{proof}
取 $D'^{[0,1]}$ 为
$[0,1]\times D'$ 的模型。
由引理 5.2.31，
存在从 $\varphi$ 到 $\psi$ 的同伦
\[
\Phi=\{\Phi_{i,j}\}_{i,j=0,1,\ldots}:
D\longrightarrow D'^{[0,1]}.
\]
由于
$D'^{[0,1]}=D'\oplus D'[-1]\oplus D'$，
分量分解可写为
\translatornote{原书实体页 298 在
$D'^{[0,1]}$ 的直和分解中
三个分量均漏去撇号，
写作 $D\oplus D[-1]\oplus D$；
这里的模型以 $D'$ 为端点对象。
同一证明的下一式还把
$\varphi_{0,0}$ 写成
$\phi_{0,0}$。
译文按已定义的对象和同态记号统一；
见疑误 A-0281。}：
\[
\Phi_{0,0}=
\varphi_{0,0}\oplus
\Phi'_{0,0}\oplus
\psi_{0,0}.
\]
由于 $\varphi\sim_s\psi$，
$\Phi_{0,0}$ 位于严格
$A_\infty$ 代数同态之上，
故
\[
\psi_{0,0}-\varphi_{0,0}
=\pm\mathfrak n'_{0,0}\circ\Phi'_{0,0}
\pm\Phi'_{0,0}\circ\mathfrak n_{0,0}.
\]
于是 (1) 成立，
(2) 是 (1) 的直接推论。
\end{proof}

\begin{remark}{5.2.39}\label{rem:5-2-39}
可以找到与式 (4.2.39) 类似的公式，
留给读者推导。
\end{remark}

对第 3.7.4--3.7.5 小节为
相对 spin 拉格朗日子流形对
$(L^{(1)},L^{(0)})$
构造的滤过 $A_\infty$ 双模
$C(L^{(1)},L^{(0)};\Lambda_{0,\mathrm{nov}})$
应用命题 5.2.37 与 5.2.38，
立即得到下述结论。

\begin{corollary}{5.2.40}\label{cor:5-2-40}
\begin{enumerate}[label=\textup{(\arabic*)},leftmargin=2.8em]
\item 若
$b_i\sim b'_i
\in\widehat{\mathcal M}(L^{(i)})$
（$i=0,1$），则
\[
\begin{aligned}
&HF\bigl((L^{(1)},b_1);
           (L^{(0)},b_0);
           \Lambda_{0,\mathrm{nov}}\bigr)\\
&\qquad\cong
HF\bigl((L^{(1)},b'_1);
          (L^{(0)},b'_0);
          \Lambda_{0,\mathrm{nov}}\bigr).
\end{aligned}
\]
\item 若
$b_i\sim b'_i
\in\widehat{\mathcal M}_{\mathrm{weak}}
(L^{(i)})$
且
$\mathfrak{PO}_{L^{(1)}}(b_1)
=\mathfrak{PO}_{L^{(0)}}(b_0)$，
则
\[
\begin{aligned}
&HF\bigl((L^{(1)},b_1);
           (L^{(0)},b_0);
           \Lambda_{0,\mathrm{nov}}\bigr)\\
&\qquad\cong
HF\bigl((L^{(1)},b'_1);
          (L^{(0)},b'_0);
          \Lambda_{0,\mathrm{nov}}\bigr).
\end{aligned}
\]
\item 若
$(\mathfrak b,b_i)\sim
(\mathfrak b',b'_i)
\in\widehat{\mathcal M}_{\mathrm{weak,def}}
(L^{(i)})$
且
$\mathfrak{PO}_{L^{(0)}}(\mathfrak b,b_0)
=\mathfrak{PO}_{L^{(1)}}(\mathfrak b,b_1)$，
则
\[
\begin{aligned}
&HF\bigl((L^{(1)},(\mathfrak b,b_1)),
          (L^{(0)},(\mathfrak b,b_0));
          \Lambda_{0,\mathrm{nov}}\bigr)\\
&\qquad\cong
HF\bigl((L^{(1)},(\mathfrak b',b'_1)),
          (L^{(0)},(\mathfrak b',b'_0));
          \Lambda_{0,\mathrm{nov}}\bigr).
\end{aligned}
\]
\end{enumerate}
上述同构均为典范同构。
\end{corollary}

第 5.2.2 小节定义了滤过
$A_\infty$ 双模结构的拉回。
拉回运算与同伦等价如下相容。

\begin{lemma}{5.2.41}\label{lem:5-2-41}
若 $\mathfrak f^{(i)}:C_i\to C'_i$
分别同伦于
$\mathfrak f^{(i)'}:C_i\to C'_i$，
且 $D$ 是 $(C'_1,C'_0)$ 上的滤过
$A_\infty$ 双模，则
\[
(\mathfrak f^{(1)},\mathfrak f^{(0)})^*D
\sim
(\mathfrak f^{(1)'},\mathfrak f^{(0)'})^*D.
\]
非滤过或弱滤过的版本，
以及／或者含幺或同伦含幺的版本，
也都成立。
\end{lemma}

\begin{proof}
设
$\mathfrak h^{(i)}:C_i\to\mathfrak C'_i$
是从 $\mathfrak f^{(i)}$ 到
$\mathfrak f^{(i)'}$ 的同伦，
其中 $\mathfrak C'_i$
是 $[0,1]\times C'_i$ 的模型。
取 $(\mathfrak C'_1,\mathfrak C'_0)$
上的滤过 $A_\infty$ 双模
$\mathfrak D$，
作为 $[0,1]\times D$ 的模型。
由引理 5.2.9，
有滤过 $A_\infty$ 双模同态
\[
\begin{aligned}
\mathrm{Eval}_{s=0}:
(\mathfrak h^{(1)},\mathfrak h^{(0)})^*\mathfrak D
&\longrightarrow
(\mathfrak f^{(1)},\mathfrak f^{(0)})^*D,\\
\mathrm{Eval}_{s=1}:
(\mathfrak h^{(1)},\mathfrak h^{(0)})^*\mathfrak D
&\longrightarrow
(\mathfrak f^{(1)'},\mathfrak f^{(0)'})^*D.
\end{aligned}
\]
由定理 5.2.35，二者均为同伦等价，
故得所求结论。
注意，上述论证不直接适用于弱滤过版本；
不过，滤过版本蕴含弱滤过版本。
\end{proof}

在各种情形下，尤其是第 5.3、5.4 节，
需要把给定的滤过 $A_\infty$ 代数
替换为与之同伦等价的其他代数。
若 $C_i$ 与 $C'_i$ 同伦等价，
还需把给定的滤过
$(C_1,C_0)$-$A_\infty$ 双模
替换为滤过
$(C'_1,C'_0)$-$A_\infty$ 双模。
这里将验证这种替换确实可行。
为清楚表述相关断言，
引入滤过 $(C_1,C_0)$-$A_\infty$
双模的范畴。

\begin{definition}{5.2.42}\label{def:5-2-42}
\begin{enumerate}[label=\textup{(\arabic*)},leftmargin=2.8em]
\item 范畴
$\mathfrak A_\infty(C_1,C_0)$
的对象是 $(C_1,C_0)$ 上的滤过
$A_\infty$ 双模 $D$。
从 $D$ 到 $D'$ 的态射是位于恒等同态之上的
滤过 $A_\infty$ 双模同态。
\item 设 $\varphi,\psi:D\to D'$ 为两个态射。
若 $\varphi\sim_{\mathrm{Incl}}\psi$，
则称 $\varphi\sim\psi$。
换言之，若存在位于
$\mathrm{Incl}:C_i\to\mathfrak C_i$
之上的滤过 $A_\infty$ 双模同态
$\Psi:D\to\mathfrak D'$，满足
\[
\mathrm{Eval}_{s=0}\circ\Psi=\varphi,
\qquad
\mathrm{Eval}_{s=1}\circ\Psi=\psi,
\]
则 $\varphi\sim\psi$。
这里 $\mathfrak C_i,\mathfrak D'$
分别是 $[0,1]\times C_i$、
$[0,1]\times D'$ 的模型。
容易看出，$\sim$ 是等价关系。
\item 定义范畴
$\mathfrak{HA}_\infty(C_1,C_0)$：
其对象与
$\mathfrak A_\infty(C_1,C_0)$
相同，态射是后者态射的
$\sim$-等价类。
\item 非滤过或弱滤过的版本，
以及／或者含幺或同伦含幺的版本，
也以同样方式定义，
分别记作
$\mathfrak{HA}_\infty(\overline C_1,\overline C_0)$、
$\mathfrak{HA}_{\infty u}(C_1,C_0)$、
$\mathfrak{HA}_{\infty hu}
(\overline C_1,\overline C_0)$ 等。
\end{enumerate}
\end{definition}

\begin{proposition}{5.2.43}\label{prop:5-2-43}
若 $\mathfrak f^{(i)}:C_i\to C'_i$
（$i=0,1$）均为同伦等价，
则拉回运算
$D\mapsto
(\mathfrak f^{(1)},\mathfrak f^{(0)})^*D$
诱导范畴等价
\[
\mathfrak{HA}_\infty(C'_1,C'_0)
\longrightarrow
\mathfrak{HA}_\infty(C_1,C_0).
\]
非滤过或弱滤过的版本，
以及／或者含幺或同伦含幺的版本，
也成立。
\end{proposition}

\begin{proof}
设 $\varphi:D\to D'$ 是
$\mathfrak A_\infty(C'_1,C'_0)$
中的态射。
由引理 5.2.9，
它诱导
$\mathfrak A_\infty(C_1,C_0)$
中的态射
\[
(\mathfrak f^{(1)},\mathfrak f^{(0)})^*\varphi:
(\mathfrak f^{(1)},\mathfrak f^{(0)})^*D
\longrightarrow
(\mathfrak f^{(1)},\mathfrak f^{(0)})^*D'.
\]
若 $\varphi\sim\psi$，
取从 $\varphi$ 到 $\psi$ 的同伦
$\mathfrak h:D\to\mathfrak D'$。
则有
\[
(\mathfrak f^{(1)},\mathfrak f^{(0)})^*\mathfrak h:
(\mathfrak f^{(1)},\mathfrak f^{(0)})^*D
\longrightarrow
(\mathfrak f^{(1)},\mathfrak f^{(0)})^*\mathfrak D'.
\]
容易看出，
$(\mathfrak f^{(1)},\mathfrak f^{(0)})^*
\mathfrak D'$
是
$[0,1]\times
(\mathfrak f^{(1)},\mathfrak f^{(0)})^*D'$
的模型。
故该拉回同伦连接
$(\mathfrak f^{(1)},\mathfrak f^{(0)})^*\varphi$
与
$(\mathfrak f^{(1)},\mathfrak f^{(0)})^*\psi$。
由此得到函子
\[
\mathfrak{HA}_\infty(C'_1,C'_0)
\longrightarrow
\mathfrak{HA}_\infty(C_1,C_0).
\]
若 $(\mathfrak f^{(i)})^{-1}$
为 $\mathfrak f^{(i)}$ 的同伦逆，
则由引理 5.2.41，
沿
$\bigl((\mathfrak f^{(1)})^{-1},
(\mathfrak f^{(0)})^{-1}\bigr)$
的拉回函子是沿
$(\mathfrak f^{(1)},\mathfrak f^{(0)})$
拉回函子的逆。
\end{proof}
